\documentclass[11pt,a4paper]{article}

\usepackage[utf8]{inputenc}
\usepackage[T1]{fontenc}
\usepackage[spanish,es-nodecimaldot]{babel}
\usepackage{geometry}
\usepackage{xcolor}
\usepackage{hyperref}
\usepackage{listings}
\usepackage{fancyhdr}
\usepackage{tabularx}
\usepackage{array}
\usepackage{longtable}
\usepackage{booktabs}
\usepackage{parskip}
\usepackage{amsmath}

\geometry{top=2.2cm,bottom=2.2cm,left=2.25cm,right=2.25cm}

\definecolor{Primary}{HTML}{174A7E}
\definecolor{Secondary}{HTML}{2A7F9E}
\definecolor{Accent}{HTML}{E07A2D}
\definecolor{LightBlue}{HTML}{EAF3F8}
\definecolor{LightGray}{HTML}{F4F6F7}
\definecolor{DarkGray}{HTML}{263238}
\definecolor{CodeGreen}{HTML}{26734D}
\definecolor{CodePurple}{HTML}{7A3E9D}

\hypersetup{
  colorlinks=true,
  linkcolor=Primary,
  urlcolor=Secondary,
  citecolor=Primary,
  pdftitle={Explicación del código de detección de objetos con YOLOv8},
  pdfsubject={Guía técnica de un notebook de detección de objetos},
  pdfauthor={Documento generado para el proyecto}
}

\pagestyle{fancy}
\fancyhf{}
\fancyhead[L]{\small\textcolor{Primary}{Detección de objetos con YOLOv8}}
\fancyhead[R]{\small\textcolor{DarkGray}{Explicación del código}}
\fancyfoot[C]{\small\thepage}
\renewcommand{\headrulewidth}{0.4pt}

\setlength{\parindent}{0pt}
\setlength{\parskip}{6pt}
\renewcommand{\arraystretch}{1.25}
\setcounter{tocdepth}{2}
\makeatletter
\renewcommand\l@section{\@dottedtocline{1}{0em}{2.8em}}
\renewcommand\l@subsection{\@dottedtocline{2}{2.8em}{3.5em}}
\makeatother

\lstdefinestyle{pythonstyle}{
  language=Python,
  basicstyle=\ttfamily\small,
  keywordstyle=\color{Primary}\bfseries,
  commentstyle=\color{CodeGreen},
  stringstyle=\color{CodePurple},
  numbers=left,
  numberstyle=\tiny\color{gray},
  stepnumber=1,
  numbersep=8pt,
  frame=single,
  rulecolor=\color{LightBlue},
  backgroundcolor=\color{LightGray},
  showstringspaces=false,
  breaklines=true,
  breakatwhitespace=false,
  columns=fullflexible,
  keepspaces=true,
  tabsize=4,
  captionpos=b,
  literate=
    {á}{{\'a}}1 {é}{{\'e}}1 {í}{{\'i}}1 {ó}{{\'o}}1 {ú}{{\'u}}1
    {Á}{{\'A}}1 {É}{{\'E}}1 {Í}{{\'I}}1 {Ó}{{\'O}}1 {Ú}{{\'U}}1
    {ñ}{{\~n}}1 {Ñ}{{\~N}}1 {ü}{{\"u}}1 {Ü}{{\"U}}1
    {¿}{{?`}}1 {¡}{{!`}}1 {º}{{\textordmasculine}}1
}

\lstset{style=pythonstyle}

\newcommand{\archivo}[1]{\nolinkurl{#1}}
\newcommand{\concepto}[1]{\textcolor{Primary}{\textbf{#1}}}
\newcommand{\nota}[1]{%
  \begin{center}
  \fcolorbox{Secondary}{LightBlue}{%
    \parbox{0.92\textwidth}{\textbf{\textcolor{Primary}{Nota.}} #1}}
  \end{center}
}
\newcommand{\advertencia}[1]{%
  \begin{center}
  \fcolorbox{Accent}{LightGray}{%
    \parbox{0.92\textwidth}{\textbf{\textcolor{Accent}{Atención.}} #1}}
  \end{center}
}

\begin{document}
\hypersetup{pageanchor=false}

\begin{titlepage}
  \pagecolor{Primary}
  \color{white}
  \vspace*{2.1cm}

  {\Huge\bfseries Explicación del código\par}
  \vspace{0.45cm}
  {\LARGE Detección de objetos en tiempo real con YOLOv8\par}

  \vspace{1.2cm}
  {\large Guía técnica y conceptual del notebook\par}
  \vspace{0.25cm}
  {\large\ttfamily NombreApellido\_DeteccionObjetos.ipynb\par}

  \vfill

  \fcolorbox{white}{Primary}{%
    \parbox{0.83\textwidth}{%
      \color{white}
      \textbf{Tecnologías principales}\\[4pt]
      Python \(\cdot\) Ultralytics YOLOv8 Nano \(\cdot\)
      OpenCV \(\cdot\) PyTorch \(\cdot\) Pandas \(\cdot\) Jupyter
    }}

  \vfill
  {\large 28 de julio de 2026\par}
  \vspace{0.6cm}
\end{titlepage}
\nopagecolor
\color{DarkGray}
\hypersetup{pageanchor=true}
\pagenumbering{roman}

\tableofcontents
\newpage
\pagenumbering{arabic}

\section{Propósito del programa}

El notebook implementa un sistema local de \concepto{visión artificial} que toma
video de la cámara de la computadora, analiza cada fotograma con el modelo
\concepto{YOLOv8 Nano} y muestra los objetos encontrados. Para cada detección
presenta tres datos:

\begin{itemize}
  \item la clase u objeto reconocido, por ejemplo \texttt{person},
  \texttt{bottle} o \texttt{cell phone};
  \item la confianza estimada por el modelo;
  \item un recuadro delimitador que indica la posición del objeto en la imagen.
\end{itemize}

El programa también permite capturar una evidencia con la tecla \textbf{C}.
La captura se conserva en memoria, se muestra dentro del notebook y se acompaña
de una tabla con las detecciones. Las teclas \textbf{Q} y \textbf{Esc} cierran
la prueba sin guardar evidencia.

\subsection{Archivos del proyecto}

\begin{tabularx}{\textwidth}{>{\raggedright\arraybackslash}p{0.38\textwidth} X}
\toprule
\textbf{Archivo o carpeta} & \textbf{Función} \\
\midrule
\archivo{NombreApellido_DeteccionObjetos.ipynb} &
Notebook que contiene las explicaciones, instalación, funciones, pruebas y
verificación final. \\
\archivo{yolov8n.pt} &
Pesos preentrenados de YOLOv8 Nano. Contienen los parámetros que el modelo
aprendió antes de usarse en este proyecto. \\
\archivo{.yolo_config} &
Carpeta local destinada a la configuración de Ultralytics, creada para evitar
problemas de permisos en carpetas globales. \\
\bottomrule
\end{tabularx}

\section{Tecnologías y responsabilidades}

\begin{longtable}{>{\raggedright\arraybackslash}p{0.23\textwidth}
                  >{\raggedright\arraybackslash}p{0.69\textwidth}}
\toprule
\textbf{Componente} & \textbf{Responsabilidad en el código} \\
\midrule
\endhead
Python & Coordina el flujo completo y conecta las bibliotecas. \\
Jupyter/IPython & Ejecuta el programa por celdas y muestra imágenes, tablas y
texto enriquecido en la salida. \\
Ultralytics YOLO & Carga el modelo y ejecuta la inferencia sobre cada
fotograma. \\
PyTorch & Proporciona el motor numérico usado por YOLO y permite comprobar si
hay una GPU CUDA disponible. \\
OpenCV (\texttt{cv2}) & Abre la cámara, obtiene fotogramas, dibuja texto,
muestra la ventana y lee las teclas. \\
NumPy & Representa las imágenes como arreglos numéricos; se importa como
\texttt{np}, aunque el notebook no lo usa directamente. \\
Pandas & Convierte las detecciones en una tabla legible. \\
Pillow (\texttt{PIL}) & Convierte el arreglo RGB en una imagen que Jupyter
puede mostrar. \\
\bottomrule
\end{longtable}

\nota{YOLO significa \emph{You Only Look Once}. Es un detector de una sola
etapa: procesa la imagen y produce clases, confianzas y posiciones en una misma
inferencia. La variante Nano se eligió por ser ligera y adecuada para una
demostración en tiempo real.}

\section{Flujo general}

\begin{center}
\setlength{\fboxsep}{7pt}
\fcolorbox{Primary}{LightBlue}{\textbf{Cámara}}
\(\longrightarrow\)
\fcolorbox{Primary}{LightBlue}{\textbf{Fotograma BGR}}
\(\longrightarrow\)
\fcolorbox{Primary}{LightBlue}{\textbf{YOLOv8}}
\(\longrightarrow\)
\fcolorbox{Primary}{LightBlue}{\textbf{Resultado}}
\end{center}
\begin{center}
\(\downarrow\)
\end{center}
\begin{center}
\fcolorbox{Secondary}{LightGray}{\textbf{Ventana anotada}}
\(\longrightarrow\)
\fcolorbox{Secondary}{LightGray}{\textbf{Tecla C}}
\(\longrightarrow\)
\fcolorbox{Secondary}{LightGray}{\textbf{Imagen + tabla en Jupyter}}
\end{center}

El ciclo se repite mientras la ventana está abierta. En cada vuelta se lee una
imagen, se ejecuta el modelo, se dibujan los resultados y se consulta el
teclado. Al capturar o salir, el bucle termina y los recursos de la cámara se
liberan.

\section{Instalación de las bibliotecas}

La primera celda de código utiliza el comando especial \texttt{\%pip}.
En Jupyter, este comando instala paquetes en el entorno asociado con el kernel
actual.

\begin{lstlisting}
%pip install -q --upgrade --no-deps \
    "torch==2.11.0+cpu" "torchvision==0.26.0+cpu" \
    --index-url https://download.pytorch.org/whl/cpu

%pip uninstall -y opencv-python-headless
%pip install -q --upgrade --force-reinstall --no-deps \
    "opencv-python>=4.8,<5"
%pip install -q "ultralytics>=8.3,<9" pandas pillow
\end{lstlisting}

\subsection{Qué hace cada instrucción}

\begin{itemize}
  \item Instala versiones de \texttt{torch} y \texttt{torchvision} preparadas
  para CPU. La opción \texttt{--no-deps} evita que \texttt{pip} modifique
  dependencias adicionales durante ese paso.
  \item Desinstala \texttt{opencv-python-headless}. Esa variante está pensada
  para servidores sin interfaz gráfica y no ofrece la ventana requerida por
  \texttt{cv2.imshow}.
  \item Reinstala la variante normal de OpenCV, que sí puede abrir ventanas.
  \item Instala Ultralytics, Pandas y Pillow. Los intervalos de versiones
  limitan cambios incompatibles.
\end{itemize}

\advertencia{Después de modificar PyTorch u OpenCV conviene reiniciar el
kernel. Un proceso de Python ya iniciado puede conservar en memoria la versión
anterior de una biblioteca.}

\section{Configuración, importaciones y carga del modelo}

\subsection{Configuración local de Ultralytics}

\begin{lstlisting}
carpeta_configuracion = Path.cwd() / ".yolo_config"
carpeta_configuracion.mkdir(exist_ok=True)
os.environ.setdefault("YOLO_CONFIG_DIR", str(carpeta_configuracion))
\end{lstlisting}

\begin{enumerate}
  \item \texttt{Path.cwd()} obtiene la carpeta desde la que se ejecuta el
  notebook.
  \item El operador \texttt{/} de \texttt{Path} construye la ruta
  \archivo{.yolo_config}.
  \item \texttt{mkdir(exist\_ok=True)} crea la carpeta si no existe y no falla
  si ya estaba creada.
  \item \texttt{setdefault} define \texttt{YOLO\_CONFIG\_DIR} solamente si no
  existía una configuración previa.
\end{enumerate}

Esta preparación se realiza \textbf{antes} de importar
\texttt{ultralytics}, de modo que la biblioteca conozca la ubicación correcta
desde el inicio.

\subsection{Carga de pesos}

\begin{lstlisting}
modelo = YOLO("yolov8n.pt")
\end{lstlisting}

La variable \texttt{modelo} queda disponible globalmente para las funciones.
Como el archivo \archivo{yolov8n.pt} ya está dentro del proyecto, la carga
puede realizarse localmente. Si no existiera, Ultralytics intentaría obtenerlo
automáticamente.

Los pesos fueron preentrenados con las categorías de COCO. Por eso
\texttt{modelo.names} relaciona cada identificador numérico con un nombre de
clase y \texttt{len(modelo.names)} informa cuántas clases puede reconocer el
modelo.

\subsection{Selección del dispositivo}

\begin{lstlisting}
dispositivo = 0 if torch.cuda.is_available() else "cpu"
\end{lstlisting}

\texttt{torch.cuda.is\_available()} pregunta si PyTorch puede usar CUDA:

\begin{itemize}
  \item si responde verdadero, se usa \texttt{0}, es decir, la primera GPU;
  \item si responde falso, se usa el procesador mediante \texttt{"cpu"}.
\end{itemize}

Con las versiones CPU instaladas explícitamente en la primera celda, lo normal
es que el resultado sea \texttt{"cpu"}. La selección automática permite que
el resto del código permanezca igual en ambos casos.

\section{Función \texttt{abrir\_camara}}

\begin{lstlisting}
def abrir_camara(indice=0):
    if os.name == "nt":
        camara = cv2.VideoCapture(indice, cv2.CAP_DSHOW)
        if not camara.isOpened():
            camara.release()
            camara = cv2.VideoCapture(indice)
    else:
        camara = cv2.VideoCapture(indice)

    camara.set(cv2.CAP_PROP_FRAME_WIDTH, 1280)
    camara.set(cv2.CAP_PROP_FRAME_HEIGHT, 720)
    return camara
\end{lstlisting}

El parámetro \texttt{indice} identifica la cámara. El valor 0 suele
representar la cámara principal; una segunda cámara normalmente se prueba con
1.

En Windows se intenta primero el backend \texttt{CAP\_DSHOW}
(DirectShow). Si no abre, el objeto se libera y se repite el intento con el
backend predeterminado. En otros sistemas se utiliza directamente el método
normal de OpenCV.

Las llamadas a \texttt{camara.set} solicitan una resolución de \(1280\times
720\). Son una petición, no una garantía: el dispositivo puede utilizar la
resolución compatible más cercana. Finalmente, la función devuelve el objeto
\texttt{VideoCapture}; la comprobación definitiva se hace en la función de
detección.

\section{Función \texttt{construir\_tabla}}

\begin{lstlisting}
def construir_tabla(resultado):
    filas = []
    cajas = resultado.boxes

    if cajas is not None:
        for numero, (clase, confianza, coordenadas) in enumerate(
            zip(cajas.cls.tolist(),
                cajas.conf.tolist(),
                cajas.xyxy.tolist()),
            start=1,
        ):
            x1, y1, x2, y2 = [round(valor)
                              for valor in coordenadas]
            filas.append({
                "N.º": numero,
                "Objeto": modelo.names[int(clase)],
                "Confianza": f"{confianza * 100:.1f} %",
                "Recuadro (x1, y1, x2, y2)":
                    f"({x1}, {y1}, {x2}, {y2})",
            })

    return pd.DataFrame(filas)
\end{lstlisting}

El objeto \texttt{resultado} contiene las predicciones de un fotograma.
\texttt{resultado.boxes} reúne todas las cajas y expone:

\begin{tabularx}{\textwidth}{>{\ttfamily}p{0.19\textwidth} X}
\toprule
\textrm{\textbf{Atributo}} & \textbf{Contenido} \\
\midrule
cls & Identificador numérico de la clase de cada detección. \\
conf & Confianza entre 0 y 1 asociada con cada caja. \\
xyxy & Coordenadas \((x_1,y_1,x_2,y_2)\) en píxeles. \\
\bottomrule
\end{tabularx}

\texttt{tolist()} mueve los datos a listas normales de Python, lo que facilita
su recorrido. \texttt{zip} alinea clase, confianza y coordenadas pertenecientes
a una misma detección. \texttt{enumerate(..., start=1)} agrega una numeración
amigable para la tabla.

La confianza se multiplica por 100 y se presenta con un decimal. La clase
numérica se convierte en nombre mediante \texttt{modelo.names}. Si no hay
objetos, \texttt{filas} permanece vacía y Pandas devuelve un
\texttt{DataFrame} sin registros.

\subsection{Sistema de coordenadas}

\begin{center}
\setlength{\fboxsep}{9pt}
\fcolorbox{Primary}{LightGray}{%
\parbox{0.72\textwidth}{
\texttt{(0, 0)} se encuentra en la esquina superior izquierda.\\
\(x\) aumenta hacia la derecha y \(y\) aumenta hacia abajo.\\[4pt]
La caja va desde \((x_1,y_1)\), esquina superior izquierda, hasta
\((x_2,y_2)\), esquina inferior derecha.
}}
\end{center}

\section{Función principal: \texttt{detectar\_en\_tiempo\_real}}

Esta función coordina la captura, la inferencia, la interfaz y la evidencia.
Recibe tres parámetros:

\begin{longtable}{>{\ttfamily\raggedright\arraybackslash}p{0.27\textwidth}
                  >{\raggedright\arraybackslash}p{0.21\textwidth}
                  >{\raggedright\arraybackslash}p{0.44\textwidth}}
\toprule
\textrm{\textbf{Parámetro}} & \textbf{Predeterminado} & \textbf{Uso} \\
\midrule
\endhead
nombre\_prueba & obligatorio & Texto que identifica la evidencia. \\
confianza\_minima & 0.30 & Descarta predicciones por debajo del umbral. \\
indice\_camara & 0 & Selecciona el dispositivo de video. \\
\bottomrule
\end{longtable}

\subsection{Apertura y validación}

\begin{lstlisting}
camara = abrir_camara(indice_camara)
if not camara.isOpened():
    raise RuntimeError(
        "No fue posible abrir la cámara..."
    )
\end{lstlisting}

La función reutiliza \texttt{abrir\_camara}. Si la cámara no quedó disponible,
lanza una excepción con recomendaciones concretas. Así evita entrar en un
bucle que nunca podría recibir imágenes.

\subsection{Lectura continua}

\begin{lstlisting}
while True:
    lectura_correcta, fotograma = camara.read()
    if not lectura_correcta:
        raise RuntimeError(
            "La cámara se abrió, pero no pudo entregar imágenes."
        )
\end{lstlisting}

\texttt{camara.read()} devuelve:

\begin{itemize}
  \item un booleano que indica si la lectura fue correcta;
  \item el fotograma, representado como un arreglo NumPy.
\end{itemize}

OpenCV almacena normalmente los canales en orden BGR
(azul, verde y rojo), no en RGB. Ese detalle será importante cuando la imagen
se envíe a Jupyter.

\subsection{Inferencia con YOLO}

\begin{lstlisting}
resultado = modelo.predict(
    source=fotograma,
    conf=confianza_minima,
    imgsz=640,
    max_det=100,
    device=dispositivo,
    verbose=False,
)[0]
\end{lstlisting}

\begin{tabularx}{\textwidth}{>{\ttfamily}p{0.24\textwidth} X}
\toprule
\textrm{\textbf{Argumento}} & \textbf{Significado} \\
\midrule
source & Imagen que se analizará. \\
conf & Umbral mínimo de confianza. \\
imgsz=640 & Tamaño de entrada utilizado para la inferencia; YOLO adapta la
imagen conservando la geometría necesaria. \\
max\_det=100 & Máximo de detecciones permitidas en un fotograma. \\
device & CPU o GPU seleccionada anteriormente. \\
verbose=False & Evita imprimir un registro por cada fotograma. \\
\bottomrule
\end{tabularx}

\texttt{predict} devuelve una colección porque puede aceptar varias imágenes.
El índice \texttt{[0]} extrae el resultado correspondiente al único fotograma
enviado.

\subsection{Dibujo e interfaz}

\begin{lstlisting}
fotograma_anotado = resultado.plot(
    line_width=2,
    font_size=13,
)

cv2.rectangle(
    fotograma_anotado,
    (0, 0),
    (fotograma_anotado.shape[1], 42),
    (25, 25, 25),
    -1,
)
cv2.putText(
    fotograma_anotado,
    "C = capturar evidencia    Q o Esc = salir",
    (12, 29),
    cv2.FONT_HERSHEY_SIMPLEX,
    0.72,
    (255, 255, 255),
    2,
    cv2.LINE_AA,
)
\end{lstlisting}

\texttt{resultado.plot()} crea una copia visual con cajas, etiquetas y
confianzas. Después, \texttt{cv2.rectangle} dibuja una franja oscura en la
parte superior. El grosor \texttt{-1} indica que el rectángulo se rellena.
\texttt{cv2.putText} coloca instrucciones blancas y
\texttt{cv2.LINE\_AA} suaviza los bordes de las letras.

El ancho de la franja se obtiene con
\texttt{fotograma\_anotado.shape[1]}, por lo que se adapta al tamaño real del
video.

\subsection{Lectura del teclado}

\begin{lstlisting}
cv2.imshow(ventana, fotograma_anotado)
tecla = cv2.waitKey(1) & 0xFF

if tecla in (ord("c"), ord("C")):
    evidencia = {
        "nombre": nombre_prueba,
        "imagen_bgr": fotograma_anotado.copy(),
        "resultado": resultado,
    }
    break

if tecla in (ord("q"), ord("Q"), 27):
    break
\end{lstlisting}

\texttt{imshow} actualiza la ventana. \texttt{waitKey(1)} concede a OpenCV un
instante para procesar eventos y obtiene la tecla. La operación
\texttt{\& 0xFF} normaliza el código a sus ocho bits inferiores.

\texttt{ord} convierte una letra en su código numérico y el número 27
representa \textbf{Esc}. Al capturar se usa
\texttt{fotograma\_anotado.copy()}; la copia impide que un cambio posterior
del arreglo altere la evidencia guardada.

\subsection{Liberación segura de recursos}

\begin{lstlisting}
finally:
    camara.release()
    cv2.destroyAllWindows()
    cv2.waitKey(1)
\end{lstlisting}

El bloque \texttt{finally} se ejecuta tanto al terminar normalmente como al
producirse un error. \texttt{release} devuelve la cámara al sistema y
\texttt{destroyAllWindows} cierra las ventanas de OpenCV. La última llamada a
\texttt{waitKey} ayuda a procesar el evento de cierre en algunos entornos.

Este diseño evita que la cámara quede ocupada y bloquee una prueba posterior.

\subsection{Presentación en Jupyter}

\begin{lstlisting}
imagen_rgb = cv2.cvtColor(
    evidencia["imagen_bgr"],
    cv2.COLOR_BGR2RGB,
)
tabla = construir_tabla(evidencia["resultado"])

display(Markdown(f"### Evidencia: {nombre_prueba}"))
display(Image.fromarray(imagen_rgb))
\end{lstlisting}

OpenCV usa BGR, mientras que Pillow y Jupyter esperan RGB.
\texttt{cvtColor} intercambia correctamente los canales para evitar que rojo
y azul aparezcan invertidos.

Luego se crea la tabla. Si está vacía, se proponen ajustes de iluminación o
confianza. Si tiene registros, se muestra sin el índice automático de Pandas y
se informa la cantidad de objetos.

La función finalmente agrega la tabla al diccionario y lo devuelve:

\begin{lstlisting}
evidencia["tabla"] = tabla
return evidencia
\end{lstlisting}

La estructura resultante contiene:

\begin{itemize}
  \item \texttt{nombre}: descripción de la prueba;
  \item \texttt{imagen\_bgr}: fotograma anotado;
  \item \texttt{resultado}: objeto original producido por Ultralytics;
  \item \texttt{tabla}: resumen tabular de las cajas.
\end{itemize}

\section{Las tres pruebas}

Las tres llamadas reutilizan la misma función y solo cambian el nombre de la
prueba y, en un caso, el umbral.

\begin{lstlisting}
evidencia_prueba_1 = detectar_en_tiempo_real(
    "Prueba 1 - Persona frente a la cámara",
    confianza_minima=0.30,
)

evidencia_prueba_2 = detectar_en_tiempo_real(
    "Prueba 2 - Objeto individual",
    confianza_minima=0.30,
)

evidencia_prueba_3 = detectar_en_tiempo_real(
    "Prueba 3 - Dos o más objetos simultáneos",
    confianza_minima=0.25,
)
\end{lstlisting}

\begin{tabularx}{\textwidth}{p{0.17\textwidth} p{0.20\textwidth} X}
\toprule
\textbf{Prueba} & \textbf{Umbral} & \textbf{Objetivo} \\
\midrule
1 & 30\,\% & Capturar una persona frente a la cámara. COCO reconoce
\texttt{person}, no el rostro como clase separada. \\
2 & 30\,\% & Detectar un objeto individual perteneciente a las clases
conocidas por el modelo. \\
3 & 25\,\% & Capturar al menos dos objetos. El umbral menor aumenta la
posibilidad de detectar objetos más difíciles. \\
\bottomrule
\end{tabularx}

Un umbral bajo aumenta la sensibilidad, pero también puede introducir falsos
positivos. Un umbral alto exige mayor seguridad al modelo, aunque puede omitir
objetos reales.

\section{Verificación final}

\begin{lstlisting}
nombres_variables = [
    "evidencia_prueba_1",
    "evidencia_prueba_2",
    "evidencia_prueba_3",
]

evidencias = [globals().get(nombre)
              for nombre in nombres_variables]
\end{lstlisting}

\texttt{globals()} devuelve el espacio de nombres global del notebook.
\texttt{get(nombre)} recupera cada variable sin provocar un
\texttt{NameError}; si una prueba no se ejecutó, devuelve \texttt{None}.

\begin{lstlisting}
if all(evidencia is not None for evidencia in evidencias):
    cantidades = [len(evidencia["tabla"])
                  for evidencia in evidencias]

    if (cantidades[0] >= 1
            and cantidades[1] >= 1
            and cantidades[2] >= 2):
        print("Se cumplen las detecciones mínimas.")
\end{lstlisting}

\texttt{all} confirma que las tres capturas existen. Después,
\texttt{len(evidencia["tabla"])} cuenta las filas detectadas. Se requieren una,
una y dos detecciones respectivamente.

\advertencia{La verificación cuenta objetos, pero no valida sus clases. Por
ejemplo, la primera prueba podría aprobar con una botella aunque no aparezca
\texttt{person}. Una validación más estricta debería comprobar también el
contenido de la columna \texttt{Objeto}.}

\section{Parámetros que modifican el comportamiento}

\begin{longtable}{>{\ttfamily\raggedright\arraybackslash}p{0.24\textwidth}
                  >{\raggedright\arraybackslash}p{0.25\textwidth}
                  >{\raggedright\arraybackslash}p{0.43\textwidth}}
\toprule
\textrm{\textbf{Parámetro}} & \textbf{Si disminuye} & \textbf{Si aumenta} \\
\midrule
\endhead
confianza\_minima &
Detecta más candidatos, incluidos posibles errores. &
Muestra menos resultados, pero generalmente más seguros. \\
imgsz &
Puede aumentar la velocidad y perder objetos pequeños. &
Puede mejorar detalle, con mayor costo de tiempo y memoria. \\
max\_det &
Limita escenas muy cargadas. &
Permite más cajas y requiere más procesamiento posterior. \\
indice\_camara &
No representa calidad; selecciona otro dispositivo. &
Prueba cámaras enumeradas con otros índices. \\
CAP\_PROP\_FRAME\_WIDTH y HEIGHT &
Reduce los datos capturados. &
Solicita más resolución, si la cámara la admite. \\
\bottomrule
\end{longtable}

\section{Manejo de errores y robustez}

El código considera varias situaciones:

\begin{itemize}
  \item intenta un backend alternativo si DirectShow no abre;
  \item verifica por separado que la cámara se haya abierto y que entregue
  fotogramas;
  \item usa \texttt{try/finally} para liberar el dispositivo;
  \item acepta mayúsculas y minúsculas para las teclas;
  \item devuelve \texttt{None} cuando el usuario sale sin capturar;
  \item maneja una tabla vacía sin asumir que hubo detecciones;
  \item recupera evidencias con \texttt{globals().get} para que la verificación
  no falle si falta una variable.
\end{itemize}

Los mensajes de solución de problemas del notebook orientan sobre permisos,
aplicaciones que ocupan la cámara, iluminación, índice de cámara, velocidad y
reinicio del kernel.

\section{Limitaciones y mejoras recomendadas}

\subsection{Limitaciones actuales}

\begin{enumerate}
  \item \textbf{Las evidencias no se exportan como archivos.} Permanecen en
  memoria y en las salidas del notebook; es necesario guardar el notebook para
  conservarlas.
  \item \textbf{La validación final solo cuenta cajas.} No comprueba que la
  primera escena contenga realmente una persona ni que la segunda corresponda
  al objeto solicitado.
  \item \textbf{Las clases aparecen en inglés.} Son los nombres originales de
  COCO proporcionados por el modelo.
  \item \textbf{Solo se usa la primera GPU.} El valor 0 no contempla selección
  entre varias GPU.
  \item \textbf{NumPy se importa pero no se referencia directamente.} Las
  imágenes sí son arreglos NumPy, pero el alias \texttt{np} podría eliminarse
  sin cambiar el comportamiento actual.
  \item \textbf{La interfaz requiere escritorio local.}
  \texttt{cv2.imshow} no es adecuada para un servidor sin pantalla o un
  notebook remoto.
\end{enumerate}

\subsection{Mejoras posibles}

\begin{itemize}
  \item Guardar automáticamente cada imagen con
  \texttt{cv2.imwrite} y exportar la tabla a CSV.
  \item Validar clases específicas, por ejemplo exigir
  \texttt{"person"} en la primera tabla.
  \item Traducir los nombres mediante un diccionario español.
  \item Calcular fotogramas por segundo para medir rendimiento.
  \item Permitir seleccionar resolución, dispositivo y umbral desde controles
  del notebook.
  \item Agregar fecha y hora a cada evidencia.
  \item Separar la lógica de inferencia de la interfaz para facilitar pruebas
  automáticas.
\end{itemize}

\section{Cómo ejecutar correctamente el notebook}

\begin{enumerate}
  \item Abrir \archivo{NombreApellido_DeteccionObjetos.ipynb} en Jupyter,
  JupyterLab o VS Code.
  \item Ejecutar la instalación una vez y reiniciar el kernel.
  \item Ejecutar la celda de importaciones y confirmar que el modelo se cargó.
  \item Ejecutar la celda que define las tres funciones.
  \item Ejecutar cada prueba por separado, dar foco a la ventana y presionar
  \textbf{C} cuando la detección sea visible.
  \item Comprobar que la imagen y la tabla aparecen debajo de cada celda.
  \item Ejecutar la verificación final.
  \item Guardar el notebook con todas sus salidas.
  \item Antes de entregar, sustituir \texttt{NombreApellido} en el archivo y
  en el encabezado por los datos reales del estudiante.
\end{enumerate}

\section{Resumen conceptual}

El programa sigue una separación de responsabilidades clara:

\begin{itemize}
  \item \texttt{abrir\_camara} se ocupa del dispositivo físico;
  \item \texttt{construir\_tabla} transforma resultados del modelo en datos
  comprensibles;
  \item \texttt{detectar\_en\_tiempo\_real} controla el ciclo completo y la
  interacción;
  \item las tres celdas de prueba reutilizan esa lógica;
  \item la última celda comprueba que las evidencias mínimas existen.
\end{itemize}

La secuencia esencial puede resumirse así:

\[
\text{capturar imagen}
\rightarrow
\text{inferir con YOLO}
\rightarrow
\text{dibujar detecciones}
\rightarrow
\text{mostrar}
\rightarrow
\text{capturar evidencia}
\]

El resultado es una aplicación educativa, local y reutilizable que conecta una
cámara con un modelo de aprendizaje profundo y convierte sus predicciones en
una salida visual y tabular.

\section*{Glosario breve}
\addcontentsline{toc}{section}{Glosario breve}

\begin{description}
  \item[Clase:] categoría predicha para un objeto.
  \item[Confianza:] puntuación con la que el modelo respalda una detección.
  \item[Fotograma:] una imagen individual de la secuencia de video.
  \item[Inferencia:] uso de un modelo entrenado para producir predicciones.
  \item[Peso:] parámetro aprendido por la red neuronal y almacenado en el
  archivo \archivo{.pt}.
  \item[Recuadro delimitador:] rectángulo que localiza un objeto.
  \item[Umbral:] valor mínimo aceptado para conservar una predicción.
\end{description}

\end{document}
